%----------------------------------------------------------------------
\section{Implementation details}
\label{app:impl}
%----------------------------------------------------------------------

\paragraph{Rendering.} Frames are $224\times224$ grayscale, replicated to three channels. Fill level
$0.12$, background $0.92$, grid lines every 28 rows at $0.80$. Values are standardised with the
statistics of the visible context, clipped to $z\in[-3,3]$ and mapped to height by
Equation~\ref{eq:height}. A clip holds $16kP$ observations in $F$ frames; tubelet size is 2, giving
$16\times14\times14 = 3136$ tokens at $F=32$.

\paragraph{Training.} VideoMAE-B initialised from Kinetics-400, \texttt{norm\_pix\_loss} off, AdamW
($\beta_1{=}0.9$, $\beta_2{=}0.95$), learning rate $10^{-4}$ with cosine decay and 500 warm-up
steps, weight decay $0.05$ (or $0$ with prior anchoring at $\lambda_{\mathrm{sp}}=0.05$), bf16,
micro-batch 8 with 4-step gradient accumulation for an effective batch of 32. Masked future length
$h_p$ is sampled from $\{2,4\}$; with probability $0.3$ the leading tubelets are additionally
hidden; $30\%$ of batches use tube masking instead. Frame scale $k\in\{1,2,4\}$ with probability
$0.5/0.3/0.2$. Learning rate $2\times10^{-4}$ is worse and $5\times10^{-5}$ ties $10^{-4}$.

\paragraph{Corpus.} 98 LOTSA subsets, 184{,}013 series, 654.2M observations, with $20\%$ synthetic
series (sinusoids with trend, noise and regime changes) mixed in. Subsets are sampled with
probability proportional to the square root of their observation count. We verified that none of
ETTh1, ETTh2, ETTm1, ETTm2, electricity or weather is present; the nearest same-domain entries are
Monash \texttt{weather} (daily Australian stations, not the Jena 10-minute set) and
\texttt{traffic\_hourly}, which is why traffic is excluded from the evaluation.

\paragraph{Evaluation.} \texttt{--mode auto} as described in Section~\ref{sec:inference}. All 24
cells use the same checkpoint. Multi-channel datasets are forecast channel by channel. Reported MSE
and MAE are computed on standardised values over stride-1 test origins.

%----------------------------------------------------------------------
\section{Per-cell comparison}
\label{app:perhorizon}
%----------------------------------------------------------------------

\begin{table}[h]
\caption{Our model, per horizon; same checkpoint and inference rule as Table~\ref{tab:main}.
Six-dataset means in the last row.}
\label{tab:perhorizon}
\begin{center}
\setlength{\tabcolsep}{4pt}
\footnotesize
\begin{tabular}{lcccccccccc}
\toprule
\multirow{2}{*}{Dataset} & \multicolumn{2}{c}{$H{=}96$} & \multicolumn{2}{c}{$H{=}192$} & \multicolumn{2}{c}{$H{=}336$} & \multicolumn{2}{c}{$H{=}720$} & \multicolumn{2}{c}{\textbf{Avg}} \\
\cmidrule(lr){2-3} \cmidrule(lr){4-5} \cmidrule(lr){6-7} \cmidrule(lr){8-9} \cmidrule(lr){10-11}
 & MSE & MAE & MSE & MAE & MSE & MAE & MSE & MAE & MSE & MAE \\
\midrule
ETTm1   & 0.292 & 0.332 & 0.332 & 0.357 & 0.363 & 0.377 & 0.423 & 0.409 & 0.352 & 0.369 \\
ETTm2   & 0.160 & 0.245 & 0.214 & 0.283 & 0.264 & 0.318 & 0.328 & 0.362 & 0.241 & 0.302 \\
ETTh1   & 0.352 & 0.380 & 0.396 & 0.408 & 0.424 & 0.425 & 0.446 & 0.449 & 0.405 & 0.416 \\
ETTh2   & 0.266 & 0.324 & 0.334 & 0.369 & 0.361 & 0.395 & 0.379 & 0.415 & 0.335 & 0.376 \\
ECL     & 0.128 & 0.219 & 0.148 & 0.239 & 0.170 & 0.262 & 0.207 & 0.294 & 0.163 & 0.254 \\
Weather & 0.139 & 0.183 & 0.181 & 0.226 & 0.232 & 0.266 & 0.308 & 0.331 & 0.215 & 0.252 \\
\midrule
Mean    & 0.223 & 0.281 & 0.267 & 0.314 & 0.302 & 0.341 & 0.349 & 0.377 & 0.285 & 0.328 \\
\bottomrule
\end{tabular}
\end{center}
\end{table}

Two baselines publish per-horizon results on this protocol: VisionTS (Table 9 of arXiv 2408.17253v4)
and VisionTS++ (Table 4, Appendix C.2 of arXiv 2508.04379v3). Table~\ref{tab:perhorizon-vts}
compares all three on every cell.

Against VisionTS we win 18 of 24 cells on MSE and 18 of 24 on MAE. The six MSE losses are ETTh1 at
$H{=}192/336/720$, ETTh2 at $H{=}192/336$, and ETTm1 at $H{=}720$: five are mid-to-long horizons on
the hourly sets, the sixth is the longest horizon of a 15-minute set.

Against VisionTS++ base, the same size class as our model, we win 17 of 24 on MSE but only 11 of 24
on MAE, with the MAE losses concentrated on weather, where VisionTS++ wins all four horizons, and on
ETTm2. This is the per-cell form of the aggregate MAE gap noted in Section~\ref{sec:main}, and it is
the clearest place where our model is behind.

% generated by paper/make_perhorizon_table.py -- do not edit by hand
\begin{table}[h]
\caption{Per-cell comparison against the two baselines that publish per-horizon results on this protocol: VisionTS (arXiv 2408.17253v4, Table 9) and VisionTS++ base (arXiv 2508.04379v3, Table 4). \textbf{Bold}: best of the three in that column group. Counts are given in the text.}
\label{tab:perhorizon-vts}
\begin{center}
\setlength{\tabcolsep}{3.4pt}
\footnotesize
\begin{tabular}{llcccccc}
\toprule
 & & \multicolumn{3}{c}{MSE} & \multicolumn{3}{c}{MAE} \\
\cmidrule(lr){3-5} \cmidrule(lr){6-8}
Dataset & $H$ & Ours & VisionTS & VisionTS++$_\mathrm{B}$ & Ours & VisionTS & VisionTS++$_\mathrm{B}$ \\
\midrule
ETTm1 & 96 & \textbf{0.291} & 0.341 & 0.316 & \textbf{0.332} & 0.347 & 0.343 \\
 & 192 & \textbf{0.332} & 0.360 & 0.347 & \textbf{0.357} & 0.360 & 0.362 \\
 & 336 & \textbf{0.362} & 0.377 & 0.368 & 0.377 & \textbf{0.374} & 0.379 \\
 & 720 & 0.423 & 0.416 & \textbf{0.408} & 0.409 & \textbf{0.405} & \textbf{0.405} \\
\midrule
ETTm2 & 96 & \textbf{0.160} & 0.228 & 0.169 & \textbf{0.245} & 0.282 & 0.248 \\
 & 192 & \textbf{0.214} & 0.262 & 0.216 & 0.283 & 0.305 & \textbf{0.279} \\
 & 336 & 0.264 & 0.293 & \textbf{0.260} & 0.318 & 0.328 & \textbf{0.308} \\
 & 720 & \textbf{0.328} & 0.343 & 0.330 & 0.362 & 0.370 & \textbf{0.358} \\
\midrule
ETTh1 & 96 & \textbf{0.352} & 0.353 & 0.369 & \textbf{0.380} & 0.383 & 0.392 \\
 & 192 & 0.396 & \textbf{0.392} & 0.399 & \textbf{0.408} & 0.410 & 0.412 \\
 & 336 & 0.424 & \textbf{0.407} & 0.415 & 0.425 & 0.423 & \textbf{0.421} \\
 & 720 & 0.446 & \textbf{0.406} & 0.424 & 0.449 & 0.441 & \textbf{0.437} \\
\midrule
ETTh2 & 96 & \textbf{0.266} & 0.271 & 0.277 & \textbf{0.324} & 0.328 & 0.326 \\
 & 192 & 0.334 & \textbf{0.328} & 0.333 & 0.369 & 0.367 & \textbf{0.362} \\
 & 336 & 0.361 & \textbf{0.345} & 0.350 & 0.395 & \textbf{0.381} & 0.384 \\
 & 720 & 0.379 & 0.388 & \textbf{0.370} & 0.415 & 0.422 & \textbf{0.409} \\
\midrule
ECL & 96 & \textbf{0.128} & 0.177 & 0.152 & \textbf{0.219} & 0.266 & 0.237 \\
 & 192 & \textbf{0.148} & 0.188 & 0.168 & \textbf{0.239} & 0.277 & 0.252 \\
 & 336 & \textbf{0.170} & 0.207 & 0.186 & \textbf{0.262} & 0.296 & 0.269 \\
 & 720 & \textbf{0.207} & 0.256 & 0.228 & \textbf{0.294} & 0.337 & 0.303 \\
\midrule
Weather & 96 & \textbf{0.139} & 0.220 & 0.145 & 0.183 & 0.257 & \textbf{0.179} \\
 & 192 & \textbf{0.181} & 0.244 & 0.187 & 0.226 & 0.275 & \textbf{0.219} \\
 & 336 & \textbf{0.232} & 0.280 & 0.240 & 0.266 & 0.299 & \textbf{0.258} \\
 & 720 & \textbf{0.308} & 0.330 & 0.317 & 0.331 & 0.337 & \textbf{0.308} \\
\bottomrule
\end{tabular}
\end{center}
\end{table}


%----------------------------------------------------------------------
\section{Initialisation at two clip lengths}
\label{app:16f}
%----------------------------------------------------------------------

\begin{table}[h]
\caption{Video against image initialisation at two clip lengths, two seeds each where available.
Five-dataset mean MSE at $H{=}96/336/720$.}
\label{tab:16f}
\begin{center}
\small
\begin{tabular}{lll}
\toprule
Clip length & Video init & Image init \\
\midrule
16 frames & 0.257/0.350/0.383, \; 0.256/0.348/0.384 & 0.268/0.359/0.400, \; 0.265/0.353/0.400 \\
32 frames & 0.245/0.330/0.379 & 0.244/0.329/0.373, \; 0.245/0.330/0.373 \\
\bottomrule
\end{tabular}
\end{center}
\end{table}

At 16 frames the video prior is worth $3$--$4\%$; at 32 frames the two initialisations tie. Read
together with the anchoring ablation (inert at 32 frames, essential at 16), this suggests that part
of what video pretraining supplies is long-range temporal context, which a longer window supplies
directly. The matched-size comparison in Table~\ref{tab:backbone} therefore uses the 5k-step budget,
where the prior is still doing work.

%----------------------------------------------------------------------
\section{Inference configuration and cost}
\label{app:inference}
%----------------------------------------------------------------------

The rule of Section~\ref{sec:inference} maps $(P, H)$ to a configuration with no per-dataset tuning.
Table~\ref{tab:passes} lists the configuration it selects for each evaluation cell and the number of
forward passes that configuration costs. \texttt{mp} (multiperiod) predicts one block of frames with $k$
periods per frame; the \texttt{sc} suffix lists the frame scales that are averaged, each of which is
a separate pass; \texttt{rollout2} feeds the prediction back and predicts twice. The system uses one
checkpoint and one rule for all cells, but it is not single-pass everywhere.

\begin{table}[h]
\caption{Inference configuration and forward passes per evaluation cell.}
\label{tab:passes}
\begin{center}
\setlength{\tabcolsep}{5pt}
\footnotesize
\begin{tabular}{llcccc}
\toprule
Dataset & $P$ & $H{=}96$ & $H{=}192$ & $H{=}336$ & $H{=}720$ \\
\midrule
ETTh1 / ETTh2 & 24 & \texttt{sc124} (3) & \texttt{sc12} (2) & \texttt{mp} (1) & \texttt{rollout2} (2) \\
ETTm1 / ETTm2 & 96 & \texttt{mp} (1) & \texttt{mp} (1) & \texttt{mp} (1) & \texttt{mp} (1) \\
ECL & 24 & \texttt{mp} (1) & \texttt{mp} (1) & \texttt{mp} (1) & \texttt{mp} (1) \\
Weather & 144 & \texttt{mp} (1) & \texttt{mp} (1) & \texttt{mp} (1) & \texttt{mp} (1) \\
\bottomrule
\end{tabular}
\end{center}
\end{table}

A 32-frame clip costs roughly $8\times$ a 16-frame clip per pass on the largest dataset, so the
inference cost of the reported system is higher than that of a 16-frame configuration even where the
pass count is one.

%----------------------------------------------------------------------
\section{Checks on the objective-mismatch result}
\label{app:diagnosis}
%----------------------------------------------------------------------

The degradation in Figure~\ref{fig:budget}a survives the following checks. Checkpoint key sets,
the \texttt{transformers} version, the attention configuration and dataloader warnings were excluded
by re-evaluating every checkpoint under matched software. The training loss is flat across the range
($0.0092 \to 0.0089$), so the model is not overfitting the pixel objective. Per cell, the 176k-step
model is $5$--$60\%$ worse than the 20k-step model, with the largest gaps at the longest horizons
(ETTh2 at $H{=}720$: $0.429 \to 0.617$). The one quantity that moves monotonically with the
schedule is the weight norm, which shrinks by roughly $30\%$ over the $9\times$ longer run under
decoupled weight decay; this is consistent with the prior-anchoring ablation helping in the
short-clip regime, though we did not isolate it as the cause. The 60k native-objective checkpoint
was never evaluated downstream, so no value is reported for it.

%----------------------------------------------------------------------
\section{Alternatives that did not work}
\label{app:failed}
%----------------------------------------------------------------------

We list these so that the obvious questions have answers, not because any of them is a finding.

\begin{table}[h]
\caption{Attempts that did not improve on the recipe in Section~\ref{sec:method}.}
\label{tab:failed}
\begin{center}
\footnotesize
\begin{tabular}{p{0.52\textwidth}p{0.40\textwidth}}
\toprule
Attempt & Outcome \\
\midrule
$9\times$ budget under the pixel objective (176k steps) & $5$--$60\%$ worse per cell \\
Larger batch (128, 256 via accumulation) & equal or worse at every horizon \\
$4\times$ more in-domain data (buildings\_900k, LargeST, Q-Traffic, Azure) & ties the corpus we use \\
Longer masked future (8 of 16 frames) & slightly worse \\
Value-loss weight $0.3$ instead of $0.1$ & worse \\
Three-checkpoint ensemble & no gain \\
Blending with a zero-parameter statistical prior & validation selection does not transfer \\
Weight-space averaging of two seeds & better on ETT, $16\%$ worse on electricity, net worse \\
Cross-corpus weight averaging & worse than either parent \\
\bottomrule
\end{tabular}
\end{center}
\end{table}

Prediction-space averaging of two seeds is the one combination that helps, by $0.6\%$ and never
worse on any cell; it appears as a secondary row in Table~\ref{tab:main}.
